\section{Why the circuit loses on Poisson: diagnostics}
\label{sec:diagnostics}

The instruments of \S\ref{sec:xai} were recorded at seven checkpoints of every
core-matrix run. Read across all five seeds, they explain the Poisson result better
than the final error alone. Figures are in Appendix~\ref{app:gallery}; the numbers below
come from the committed \texttt{xai/*.npz} files
(\texttt{scripts/report/fig\_xai.py}).

\paragraph{The circuit learns the hard frequency and loses the easy one.} The Poisson
solution is $\sin\pi x+0.3\sin15\pi x$. Tracking the error spectrum at those two
frequencies (Figure~\ref{fig:specerr_gallery}), \texttt{q\_serial} cuts its $15\pi$
error to 2--13\% of its starting value by step 1{,}000 in four of five seeds (median
10\%, against 103\% for \texttt{c\_ff}, 86\% for \texttt{c\_mlp} and 51\% for
\texttt{c\_rff\_matched}). Its $\pi$ error falls to about 30\% by step 5{,}000 and then \emph{rises} to
1.9--2.0$\times$ its starting value by step 20{,}000, in all five seeds. The classical
model given the same two frequencies, \texttt{c\_rff\_matched}, reduces both. This is
why longer training made \texttt{q\_serial} worse, and why its seeds agree so closely
(worst over best seed $1.1\times$, Figure~\ref{fig:seedspread}): it converges reliably
to the same wrong trade-off.

\paragraph{A plausible mechanism.} The residual loss weights each Fourier mode of the
error by $\omega^2$, so the $15\pi$ mode carries $(15)^2=225$ times the weight of the
$\pi$ mode. A linear head on fixed sine features (\texttt{c\_rff\_matched}) can set the
two amplitudes independently. The circuit cannot: all 81 amplitudes of its output are
functions of the same ten angles, read out through a single expectation value and a
one-weight head. Our reading, which these data support but do not prove, is that the optimiser
buys the heavily weighted $15\pi$ amplitude at the cost of the lightly weighted $\pi$
amplitude because the two are coupled. The loss agrees: it stalls at
$\approx3.3\times10^4$ from step 1{,}000 onward while the relative error keeps moving.

\paragraph{The circuit uses few directions in parameter space.} The NTK's effective rank
at the end of training is 1.4 for \texttt{q\_serial} on Poisson and 1.1 on heat, against
6.6 and 16.9 for the 129-parameter Fourier-feature model; the Fisher effective dimension
is 3.6 and 3.7, against 18.0 and 57.4 (Figure~\ref{fig:evolution}). The size-matched
classical model \texttt{c\_rff\_matched} (13--17 parameters) is similar (2.0 and 1.9;
3.2 and 4.4), so this is mostly a matter of size, not of being quantum. The quantum
models' effective dimension barely moves during training while \texttt{c\_mlp}'s grows
from about 4 to 12 (Poisson) and from 5 to 50 (heat): the circuits never leave the few
directions they start with, consistent with PR-11's refutation.

\paragraph{The design rule's positive result is visible directly.} On heat, where
PR-6 is confirmed, the random-scaling twin \texttt{q\_random} is the least reliable
model in the whole matrix (worst over best seed $7.4\times$), while \texttt{q\_serial}'s
seeds agree to within 4\%. The designed frequencies do not make the circuit accurate,
but they make it consistently better than the same circuit with the wrong frequencies.

\paragraph{What we tested and what we did not.} A second candidate mechanism is the
hard boundary ansatz $u=B+x(1-x)N$, which asks $N$ to learn $u/(x(1-x))$, a function
that is not band-limited even when $u$ is. A post-submission pilot of an alternative
ansatz that keeps band-limited solutions representable
(\texttt{models.base.AffineBCWrapper}) gave a \emph{worse} error after 3{,}300 steps
(1.24 against 0.50), so we do not claim this mechanism. Its full test is part of the
v2 study (\S\ref{sec:followup}).
